Предметом настоящей монографии являются детерминированные
микроструктуры~-- регулярные элементы заданной геометрии (лунки,
канавки, текстуры), описываемые явным полем добавки~$\Delta h$
в суперпозиции зазора~\eqref{eq:h_superposition} и параметризованные
в подразделе~2.2. Наряду с детерминированным рельефом любая
обработанная поверхность содержит шероховатость~-- случайную
составляющую топографии, возникающую при механической обработке,
полировке или лазерном текстурировании. Шероховатость
характеризуется не явной геометрией, а статистическими параметрами:
среднеарифметическим отклонением профиля~$R_a$,
среднеквадратичным отклонением~$R_q$ и длиной
корреляции~$\lambda_c$. Настоящий раздел устанавливает, при каких
условиях шероховатость существенно влияет на решение задачи
и какие подходы позволяют учесть её при необходимости.

\textit{Критерий значимости шероховатости.}
Определяющим безразмерным параметром является отношение
характерной толщины смазочного слоя~$h_0$ к среднеквадратичной
шероховатости:
\begin{equation}
  \Lambda_r = \frac{h_0}{R_q}\,,
  \label{eq:lambda_roughness}
\end{equation}
\begin{conditions}
$h_0$ -- характерная толщина зазора, м;\\
$R_q$ -- комбинированная среднеквадратичная шероховатость
контактирующих поверхностей, м.
\end{conditions}

Параметр~$\Lambda_r$ называется числом смазочного слоя
(\textit{film thickness ratio}) и определяет режим взаимодействия
поверхностей. При $\Lambda_r \gg 1$ (условно $\Lambda_r > 5$)
поверхности надёжно разделены плёнкой смазки, высота
шероховатостей пренебрежимо мала по сравнению с толщиной зазора,
и поле давления определяется макрогеометрией и детерминированным
микрорельефом. При $\Lambda_r \approx 1\text{--}3$ высота
микронеровностей становится соизмеримой с толщиной зазора,
шероховатость начинает взаимодействовать с полем давления, и для
адекватного описания потоков требуются модифицированные уравнения.
При $\Lambda_r < 1$ реализуется режим граничной смазки
с непосредственным контактом микронеровностей; рейнольдсовская
модель тонкого слоя в этом режиме неприменима.

Для типичных подшипниковых задач с детерминированными
микроструктурами значение~$\Lambda_r$ составляет 3--10, что
соответствует режиму гидродинамической смазки. Глубина ячейки
микрорельефа~$h_p = H_p \cdot h_0$ (подраздел~2.2.4,
формула~\eqref{eq:Hp_nondim}) при типичных значениях
$H_p \sim 0{,}1\text{--}1$ существенно превышает амплитуду
шероховатости: $h_p \gg R_q$. Иерархия масштабов~-- базовый
зазор, детерминированная текстура и шероховатость~-- схематически
представлена на рисунке~\ref{fig:roughness_scales}. Это
соотношение масштабов обосновывает допущение о гладкости
поверхностей, принятое в расчётах глав~5--7.

\begin{figure}[!ht]
\centering
\fbox{\parbox{0.85\textwidth}{\centering
  \vspace{3.5cm}
  ПЛЕЙСХОЛДЕР ДЛЯ РИСУНКА\\
  Схема масштабов топографии поверхности:\\
  три уровня: базовый зазор $h_{\mathrm{base}}$,\\
  детерминированная текстура $\Delta h_{\mathrm{texture}}$ (глубина $H_p h_0$),\\
  шероховатость $\Delta h_{\mathrm{rough}}$ (амплитуда $R_q$).\\
  Показать что $R_q \ll H_p h_0$ в типичном случае.
  \vspace{3.5cm}
}}
\caption{Иерархия масштабов топографии поверхности:
  базовый зазор, детерминированная микротекстура и шероховатость}
\label{fig:roughness_scales}
\end{figure}

\textit{Подходы к учёту шероховатости в уравнении Рейнольдса.}
В литературе сложились три основных подхода, различающихся
по сложности реализации и области применимости.

Первый подход~-- детерминированный: шероховатость задаётся явным
полем $\Delta h_{\mathrm{rough}}(s_1, s_2)$, полученным по данным
измерений (атомно-силовая микроскопия, контактная профилометрия),
и добавляется в суперпозицию зазора наряду с текстурой:
\[
  h = h_{\mathrm{base}} + \Delta h_{\mathrm{texture}}
    + \Delta h_{\mathrm{rough}}.
\]
Такой подход требует измеренного профиля поверхности
и разрешения расчётной сетки, достаточного для представления
случайного рельефа: шаг сетки должен быть существенно меньше
длины корреляции~$\lambda_c$ (единицы микрометров). Это приводит
к значительному увеличению числа узлов по сравнению с задачей для
гладких поверхностей (глава~3) и обоснованно только в рамках
исследовательских расчётов.

Второй подход~-- стохастический, основанный на осреднении
уравнения Рейнольдса по ансамблю реализаций случайной
шероховатости. Пионерская модель Кристенсена (Christensen, 1969)
и её развитие в работах Патира и Ченга (Patir, Cheng, 1978--1979)
вводят поправочные коэффициенты (\textit{flow factors}), которые
учитывают влияние статистической шероховатости на средние потоки
смазки. Осреднённое уравнение Рейнольдса записывается для
среднего давления~$\bar{p}$ с модифицированными коэффициентами:
\begin{equation}
  \frac{\partial}{\partial s_1}\!\left(\phi_p\,\bar{h}^3\,
  \frac{\partial \bar{p}}{\partial s_1}\right)
  + \Lambda^2\frac{\partial}{\partial s_2}\!\left(\phi_p\,\bar{h}^3\,
  \frac{\partial \bar{p}}{\partial s_2}\right)
  = \phi_s\,\frac{\partial \bar{h}}{\partial s_1}
  + \phi_c\,\sigma_r\,\frac{\partial \phi_s}{\partial s_1}\,,
  \label{eq:reynolds_rough_avg}
\end{equation}
\begin{conditions}
$\bar{h}$ -- номинальная (без шероховатости) толщина зазора;\\
$\bar{p}$ -- осреднённое давление;\\
$\phi_p$ -- фактор давления (pressure flow factor),
корректирующий пуазейлевский поток;\\
$\phi_s$ -- фактор сдвига (shear flow factor),
корректирующий куэттовский поток;\\
$\phi_c$ -- контактный фактор;\\
$\sigma_r$ -- суммарная среднеквадратичная шероховатость
контактирующих поверхностей, м;\\
$\Lambda$ -- отношение масштабов длины
(подраздел~2.1.6).
\end{conditions}

Факторы~$\phi_p$, $\phi_s$ зависят от отношения~$\bar{h}/\sigma_r$
и от ориентации шероховатости (продольная, поперечная,
изотропная). Они вычисляются однократно по статистическим
параметрам поверхности и не требуют разрешения случайного поля
на расчётной сетке, что делает стохастический подход
вычислительно эффективным. Структура
уравнения~\eqref{eq:reynolds_rough_avg} сохраняет форму исходного
уравнения Рейнольдса~\eqref{eq:reynolds_full}, поэтому солверы
и алгоритмы, описанные в главе~3, могут быть адаптированы
с минимальными изменениями.

Третий подход~-- пренебрежение шероховатостью: при $\Lambda_r \gg 1$
полагается $\Delta h_{\mathrm{rough}} = 0$, и уравнение Рейнольдса
решается для номинального зазора~\eqref{eq:h_superposition}. Это
допущение принято в расчётах глав~5--7 настоящей монографии.

\textit{Взаимодействие шероховатости с детерминированными
микроструктурами.}
При совместном присутствии детерминированных микроструктур
и случайной шероховатости характерные масштабы могут различаться
на порядки. Глубина ячейки микрорельефа~$h_p = H_p \cdot h_0$
составляет, как правило, 0,1--1~мкм при характерном размере
ячейки в десятки микрометров; шероховатость полированных
поверхностей характеризуется значениями $R_q \sim 0{,}01\text{--}0{,}1$~мкм
при длине корреляции $\lambda_c \sim 1\text{--}10$~мкм.

Когда $R_q \ll H_p \cdot h_0$~-- типичный случай при хорошо
обработанных поверхностях~-- детерминированные микроструктуры
доминируют в формировании поля давления и потоков, а вклад
шероховатости сводится к слабому стохастическому возмущению,
не влияющему на интегральные характеристики подшипника.

Когда $R_q$ оказывается соизмерима с~$H_p \cdot h_0$, что
возможно при грубой обработке или мелком микрорельефе,
шероховатость модифицирует распределение давления внутри ячеек
текстуры, изменяет локальные кавитационные зоны
(раздел~4.1) и может как усилить, так и ослабить эффективность
текстурирования. В этом режиме результаты, полученные в рамках
модели гладких поверхностей, требуют поправок, и корректное
описание предполагает использование одного из рассмотренных выше
подходов~-- детерминированного или стохастического.

Практический вывод для проектирования состоит в том, что
глубина детерминированных микроструктур~$h_p$ должна существенно
превышать среднеквадратичную шероховатость ($H_p \cdot h_0 \gg R_q$),
чтобы регулярный рельеф не терял функциональности на фоне случайных
микронеровностей. Это условие одновременно обеспечивает
$\Lambda_r \gg 1$, при котором пренебрежение шероховатостью
в уравнении Рейнольдса обосновано.

В расчётах глав~5--7 поверхности приняты гладкими:
$\Delta h_{\mathrm{rough}} = 0$. Допущение оправдано при
$\Lambda_r \gg 1$ и $H_p \cdot h_0 \gg R_q$, что типично
для прецизионных подшипников с полированными поверхностями.
Стохастический подход с коэффициентами потоков~$\phi_p$,
$\phi_s$~\eqref{eq:reynolds_rough_avg} представляет наиболее
практичный путь расширения модели на случай существенной
шероховатости: он сохраняет структуру дискретизации и не требует
разрешения случайного поля на расчётной сетке.
Детерминированный подход применим при наличии измеренных профилей
поверхности и вычислительных ресурсов для высокого разрешения
сетки, характерного масштаба $\lambda_c$.
